\renewcommand{\bibname}{第 1 章参考文献}

\begin{thebibliography}{99}\setlength{\itemsep}{0.25em}
\bibitem{braden} H.~W.~Braden and A.~Marshakov, Singular phases of the Seiberg--Witten
integrable systems, \emph{Nucl. Phys. B} \textbf{595} (2001), 417; arXiv:hep-th/0009060.
\bibitem{cimasoni} D.~Cimasoni and N.~Reshetikhin, Dimers on surface graphs and spin
structures I, \emph{Comm. Math. Phys.} \textbf{275} (2007), 187--208; arXiv:math-ph/0608070.
\bibitem{cohen} M.~Cohen, O.~Dasbach, and H.~Russell, A twisted dimer model for knots,
\emph{Fund. Math.} \textbf{225}(1) (2014), 57--74; arXiv:1010.5228.
\bibitem{fockgoncharov1} V.~V.~Fock and A.~B.~Goncharov, Cluster $X$-varieties,
amalgamation and Poisson--Lie groups, in \emph{Algebraic Geometry Theory and Number
Theory}, pp.~27--68, Progr. Math.~253, Birkhäuser Boston, Boston, MA, 2006;
math.RT/0508408.
\bibitem{fockgoncharov2} V.~V.~Fock and A.~B.~Goncharov, Moduli spaces of local systems
and higher Teichmüller theory, \emph{Publ. Math. Inst. Hautes Études Sci.} \textbf{103}(1)
(2006), 1--211; arXiv:math/0311149.
\bibitem{fockmarshakov} V.~V.~Fock and A.~Marshakov, A note on quantum groups and
relativistic Toda theory, \emph{Nucl. Phys. B (Proc. Suppl.)} \textbf{56}(3) (1997),
208--214.
\bibitem{fomin} S.~Fomin and A.~Zelevinsky, Double Bruhat cells and total positivity,
\emph{J. Amer. Math. Soc.} \textbf{12}(2) (1999), 335--380; arXiv:math/9802056v1.
\bibitem{fulton} W.~Fulton, \emph{Introduction to Toric Varieties}, Annals of Mathematics
Studies 131, Princeton University Press, New Jersey, 1993.
\bibitem{gekhtman} M.~Gekhtman, M.~Shapiro, S.~Tabachnikov, and A.~Vainshtein, Higher
pentagram maps, weighted directed networks, and cluster dynamics, \emph{Electron. Res.
Announc. Math. Sci.} \textbf{19} (2012), 1--17; arXiv:1110.0472.
\bibitem{glick} M.~Glick, The pentagram map and $Y$-patterns, \emph{Adv. Math.}
\textbf{227}(2) (2011), 1019--1045; arXiv:1005.0598.
\bibitem{goncharovkenyon} A.~Goncharov and R.~Kenyon, Dimers and cluster integrable
systems, \emph{Ann. Sci. Éc. Norm. Supér.} \textbf{46} (2013), 747--813;
arXiv:1107.5588.
\bibitem{hoffman} T.~Hoffman, J.~Kellendonk, N.~Kuntz, and N.~Reshetikhin, Factorization
dynamics and Coxeter--Toda lattices, \emph{Comm. Math. Phys.} \textbf{212}(2) (2000),
297--321; arXiv:solv-int/9906013.
\bibitem{hofstadter} D.~R.~Hofstadter, Energy levels and wave functions of Bloch electrons
in rational and irrational magnetic fields, \emph{Phys. Rev. B} \textbf{14} (1976),
2239--2249.
\bibitem{johnson} D.~Johnson, Spin structures and quadratic forms on surfaces,
\emph{J. London Math. Soc.} \textbf{22}(2) (1980), 365--373.
\bibitem{kac} V.~Kac, \emph{Infinite Dimensional Lie Algebras} (third edition), Cambridge
University Press, Cambridge, 1990. xxii+400 pp.
\bibitem{kasteleyn} P.~W.~Kasteleyn, Graph theory and crystal physics, in F.~Harary,
editor, \emph{Graph Theory and Theoretical Physics}, Academic Press, 1967.
\bibitem{khesin1} B.~Khesin and F.~Soloviev, The pentagram map in higher dimensions and
KdV flows, \emph{Electron. Res. Announc. Math. Sci.} \textbf{19} (2012), 86--96;
arXiv:1205.3744.
\bibitem{khesin2} B.~Khesin and F.~Soloviev, Integrability of higher pentagram maps,
\emph{Math. Ann.} \textbf{357}(3) (2013), 1005--1047; arXiv:1204.0756.
\bibitem{kogan} M.~Kogan and A.~Zelevinsky, On symplectic leaves and integrable systems in
standard complex semisimple Poisson--Lie groups, \emph{Int. Math. Res. Not.} \textbf{32}
(2002), 1685--1702; arXiv:math/0203069.
\bibitem{krichever} I.~Krichever and K.~Vaninsky, The periodic and open Toda lattice,
\emph{Mirror symmetry IV} (2000), 139--158, AMS/IP Stud. Adv. Math.~33, Amer. Math.
Soc., Providence, RI, 2002; arXiv:hep-th/0010184.
\bibitem{marshakov} A.~Marshakov, Lie groups, cluster variables and integrable systems,
\emph{J. Geom. Phys.} \textbf{67} (2013), 16--36; arXiv:1207.1869.
\bibitem{ovsienko1} V.~Ovsienko, R.~Schwartz, and S.~Tabachnikov, The pentagram map: a
discrete integrable system, \emph{Comm. Math. Phys.} \textbf{299}(2) (2010), 409--446;
arXiv:0810.5605.
\bibitem{ovsienko2} V.~Ovsienko, R.~Schwartz, and S.~Tabachnikov, Liouville--Arnold
integrability of the pentagram map on closed polygons, \emph{Duke Math. J.}
\textbf{162}(12) (2013), 2149--2196; arXiv:1107.3633.
\bibitem{reiman} A.~Reiman and M.~Semenov-Tian-Shansky, {\fontspec{Times New Roman}\itshape Интегрируемые системы} 352,
{\fontspec{Times New Roman}Москва--Ижевск, РХД, 2003}.
\bibitem{ruijsenaars} S.~N.~M.~Ruijsenaars, Relativistic Toda systems,
\emph{Comm. Math. Phys.} \textbf{133} (1990), 217.
\bibitem{schwartz} R.~Schwartz, The pentagram map is recurrent, \emph{Exper. Math.}
\textbf{10} (2001), 519--528.
\bibitem{soloviev} F.~Soloviev, Integrability of the pentagram map, \emph{Duke Math. J.}
\textbf{162}(15) (2013), 2815--2853; arXiv:1106.3950.
\bibitem{suris} Y.~B.~Suris, Algebraic structure of discrete-time and relativistic Toda
lattices, \emph{Phys. Lett. A} \textbf{156} (1991), 467--474.
\bibitem{thurston} D.~Thurston, From dominoes to hexagons, arXiv:math/0405482.
\end{thebibliography}
